\section{実験計画法2；Between計画}
\label{D1_lesson13}

\subsection{授業内容}
\subsubsection{科目の中でのこのコマの位置づけ}
\begin{tabular}[t]{ccccccc}
    記述統計[4] & 確率[2] & モデル[6] & 推測・推定[1] & 意思決定[0] & 実践[2]
\end{tabular}
\subsubsection{概要}
実験計画に基づいて，データをどのように扱っていくかを具体的に確認する。まずは一要因一要因3水準，Betweenデザインの数値例を用いる。ここで改めて，標本平均と母平均の違いや平均因果効果の推論をしようとしていることを確認し，統計的に(すなわち母集団において)差があるかないかをどのように考えるかを理解する。回帰分析の時と同じように線形モデルで考え，数式の表現を行い，回帰分析と同じ形式Formであることを確認する。
\subsubsection{コマ主題細目}
\begin{description}

\item[デザイン行列]効果と誤差をそれぞれ未知母数$\delta_j,\epsilon_{ij}$とおき，実測値${y_{ij}}$がどのように表現できるかを考える。この時，効果の有無を0と1で表した説明変数$X$を用いることになる。このように数式で表現することで，その形が回帰分析と同じであること，説明変数のデータ種が違うだけであることを再確認できる。
誤差の算出の時とは異なり，回帰分析的表現では切片を第一水準の平均値におき，そこからの相対比較で表現していた。これを全体平均からの相対比較に書き換えることで，両者の数学的表現を同一にすることが可能である。ただしこの場合は，推定できる効果の未知母数が「増える」のではなく，数学的性質から「制約をかける」ことで，実質的に推定すべき母数の数は同じであることが確認できる。
\begin{flushright}
    $\to$ \citet{kosugi2018}のP.136--138
\end{flushright}

\item[二要因のモデル]二要因の実験計画法モデルを考えるために，数値例を導入する。まず数値を可視化し，次に要因ごとに周辺化することで要因の効果をどのように書き出すことができるかを考える。回帰式に書き起こすことで，重回帰分析と同じ形式Formになっていることを確認する。
\begin{flushright}
    $\to$データは\citet{Yamada2004}のP.184,表7.6.1を用いる。計算式もP.184--189を参照すると良い。
\end{flushright}

\item[効果の算出と平方和の分解]回帰式を参照しながら，平方和を分解する計算を進める。2つの要因それぞれの大きさを各要因の水準ごとに考えることで計算する。要因ごとに区切るのは周辺化と呼ばれる。可視化したグラフを参考にしながら，計算箇所を辿りながら考えることが重要である。最後に2つの効果を全体平均に加え，水準の平均と合致「しない」ことを確認する。
\begin{flushright}
    $\to$ \citet{Yamada2004}のP.184--189を参照すると良い。
\end{flushright}

\item[交互作用]合計が合致しなかったのは，組み合わせによる効果が発生していることによる。このような効果のことを交互作用interactionという。交互作用の大きさを見積もることで平方和を完全に分解できたことになる。交互作用はその出現の仕方に様々なパターンがあることを理解する。また組み合わせ爆発という言葉があるように，三要因以上の実験計画を立てると，考慮すべき交互作用項がどんどん増えていくので，実質的に研究として使える実験計画は2，3要因の計画にとどめておいた方が良いことを理解する。
\begin{flushright}
    $\to$\citet{Yamada2004}のP.190--193.\citet{kosugi2019}のP.139--142
\end{flushright}


\end{description}
\subsubsection{キーワード}
\begin{itemize}
\item 誤差と効果
\item 一般線形モデル
\item 主効果
\item 交互作用
\item 多重共線性
\end{itemize}
\subsection{授業情報}
\paragraph{コマの展開方法}
講義
\paragraph{標準シラバスにおける位置づけ}
科目番号4；心理学研究法；2.データを用いた実証的な思考方法；C データの統計的記述\\
科目番号5；心理学統計法；1.心理学で用いられる統計手法；G 推測統計:代表値と散布度をめぐって(1)
\subsubsection{予習・復習課題}
\paragraph{予習}
統計環境Rを使った線形モデルを実際に行い，切片と傾きが算出されることを確認し，また説明変数が離散的であった場合の切片や傾きがどういう意味を持っているのかを改めて考えておくと理解が進む。また，実験計画と回帰分析が同じである，という原理を十分に理解し授業に臨む必要がある。またその観点から，複数の変数で行う回帰分析，すなわち重回帰分析についても復習しておく必要がある。
\paragraph{復習}
全体平均，群平均を計算に用いて効果と誤差を算出する過程を，もう一度確認しておく。この時，自分で収集した適当なデータがあると理解が進む。身の回りにあるデータ(ex.お菓子や食べ物の重さなど)や，公開されているデータ(ex.スポーツに関するデータなどは公開されていることが多い)を使って計算するとよい。計算にはExcel/Numbers/Calcsなどの表計算ソフトを用いても良いし，統計環境Rで行列のデータを対象に分析してももちろん良い。


